% --- Άσκηση 14801 (14801.tex) ---
\Askhsh{14801}{1}
\begin{ylh}{2}
    \item 2.3. Απόλυτη Τιμή Πραγματικού Αριθμού
\end{ylh}
\begin{alist}
\item Να χαρακτηρίσετε τις προτάσεις που ακολουθούν, γράφοντας δίπλα στον αριθμό που αντιστοιχεί σε κάθε πρόταση Σ (Σωστό), αν η πρόταση είναι αληθής ή Λ (Λάθος), αν η πρόταση είναι ψευδής.

\begin{rlist}
\item Για οποιουσδήποτε πραγματικούς αριθμούς $a, \beta, \gamma, \delta$ ισχύει η πρόταση: Αν $a < \beta$ και $\gamma < \delta$, τότε $a\gamma < \beta\delta$.

\item Για κάθε $\theta \in (0, +\infty)$ ισχύει: $|x| < \theta \Leftrightarrow -\theta < x < \theta$.

\item Η εξίσωση $x^3 = 5$ έχει δύο πραγματικές ρίζες.

\item Αν ισχύουν $a > 0$ και $\Delta < 0$, όπου $\Delta$ η διακρίνουσα του τριωνύμου $a x^2 + \beta x + \gamma$, τότε το τριώνυμο $a x^2 + \beta x + \gamma$ είναι θετικό για οποιονδήποτε πραγματικό αριθμό $x$.

\item Ο παρακάτω πίνακας θα μπορούσε να είναι πίνακας τιμών μιας συνάρτησης $f$ με πεδίο ορισμού το διάστημα $[0, 4]$.
\begin{center}
\begin{tblr}{hlines = {1pt,white},
vlines = {0.5pt,white},
column{odd} = {maincolor!40,c},
column{even} = {maincolor!10,c},
column{1} = {maincolor,c,fg={white}}}
$x$ & 0 & 1 & 1 & 2 & 4 \\
$y$ & 2 & 3 & 4 & 5 & 6 \\
\end{tblr}
\end{center}
\monades{10}
\end{rlist}
\item Να αποδείξετε ότι, για οποιουσδήποτε πραγματικούς αριθμούς $a, \beta$ ισχύει η ανισότητα: $|a + \beta| \le |a| + |\beta|$. \monades{15}
\end{alist}
